\documentclass{article}
\usepackage{amsmath}
\usepackage{amssymb}
\usepackage{amsthm}
\usepackage[T1]{fontenc}


\newtheorem*{conjecture}{Conjecture}
\newtheorem*{theorem}{Theorem}
\newtheorem{lemma}{Lemma}
\newtheorem*{lemma*}{Lemma}
\newtheorem*{proposition}{Proposition}


\makeatletter
\newcommand*{\rom}[1]{\expandafter\@slowromancap\romannumeral #1@}
\makeatother



\begin{document}

\section*{Problem 1}

\subsection*{a.}
The range, \(R_{X}\) is:
\begin{equation*}
    R_{X} = \{0, 1, 2\}
\end{equation*}
\subsection*{b.}
\begin{equation*}
    P(X \geq 1.5) = 1 - P(X < 1.5) = 1 - \sum_{x \in \{0, 1\}}P_{X}(x)
\end{equation*}
\begin{equation*}
    = 1 - \Big(P(X=0) + P(X=1)\Big) = 1 - \frac{1}{2} + \frac{1}{3} = \frac{1}{6}
\end{equation*}
\subsection*{c.}
\begin{equation*}
    P(0 < X < 2) = P(X \in (0, 2)) = P(X \in \{1\}) = P(X = 1) = \frac{1}{3}
\end{equation*}
\subsection*{d.}
\begin{equation*}
    P(X=0|X<2) = \frac{P(X=0, X < 2)}{P(X < 2)} = \frac{P(X = 0)}{P(X < 2)} = 
\end{equation*}
\begin{equation*}
    = \frac{\frac{1}{2}}{P(X=0) + P(X=1)} = \frac{\frac{1}{2}}{\frac{1}{2} + \frac{1}{3}}
\end{equation*}
\begin{equation*}
    = \frac{\frac{1}{2}}{\frac{5}{6}} = \frac{3}{5}
\end{equation*}

\section*{Problem 2}

Given the data we know:

\begin{equation*}
    P(X > 3) = 0 = P(X < 0) \implies R_{X} = \{0, 1, 2, 3\}
\end{equation*}
\begin{equation*}
    P(X=2) = P(X=1)
\end{equation*}
\begin{equation*}
    P(X=0) = P(X=3)
\end{equation*}
\begin{equation*}
   P(X=1 \ \mbox{or} \ X=2) = \frac{P(X=0 \ \mbox{or} \ X=3)}{2}
\end{equation*}

Notice that by the properties of preimages (Appendix):
\begin{equation*}
    P(X=1 \ \mbox{or} \ X=2) = P(X=1) + P(X=2)
\end{equation*}
\begin{equation*}
    P(X=0 \ \mbox{or} \ X=3) = P(X=0) + P(X=3)
\end{equation*}
then:
\begin{equation*}
\begin{cases}
      P(X=1) + P(X=2) = \frac{P(X=0) + P(X=3)}{2} \\
      P(X=1) + P(X=2) + P(X=0) + P(X=3) = 1
\end{cases}  
\end{equation*}
%
Let \(A = P(X=1) + P(X=2)\) and \(B = P(X=0) + P(X=3)\). We can write the system of equations in the following way:
\begin{equation*}
    \begin{cases}
      A = \frac{B}{2} \\
      A + B = 1
\end{cases}  
\end{equation*}
Solving above gives us \(A = \frac{1}{3}\), \(B = \frac{2}{3}\). Then \(P(X=1) + P(X=2) = \frac{1}{3}\) and we immediately
get \(P(X=1) = P(X=2) = \frac{1}{6}\). By analogy, we have \(P(X=0) = P(X=3) = \frac{1}{3}\). Finally, we write the PMF:
\begin{equation*}
    P_{X}(x) = 
    \begin{cases}
        \frac{1}{3}, & \mbox{for \(x \in \{0, 3\}\)} \\
        \frac{1}{6}, & \mbox{for \(x \in \{1, 2\}\)} \\
        0, & \mbox{otherwise}
    \end{cases}
\end{equation*}
\subsubsection*{Appendix - preimages}
For \(P(X=1 \ \mbox{or} \ X=2)\) we have:
\begin{equation*}
    P(X=1 \ \mbox{or} \ X=2) = P\Bigg(\{\omega \in \Omega: X(\omega) = 1\} \cup \{\omega \in \Omega: X(\omega) = 2\}\Bigg) =
\end{equation*}
\begin{equation*}
    = P(\{\omega \in \Omega: X(\omega) \in \{1, 2\}\}) = P(X \in \{1, 2\})
\end{equation*}
\begin{equation*}
    = \sum_{x \in \{1,2\}}P_{X}(x) = P(X=1) + P(X=2)
\end{equation*}
We repeat the same for \(P(X=0 \ \mbox{or} \ X=3)\).



\section*{Problem 3}
 Let \(Z = X - Y\). The range of \(R_{Z}\) is \(\{-5, -4, \dots, 0, \dots, 5\}\). Because
\(\#\Omega = 36\) we can write down the following PMF:

\begin{equation*}
    P_{Z}(x) = 
    \begin{cases}
        \frac{6}{36} = \frac{1}{6}, & \mbox{for \(x \in \{0\}\)} \\
        \frac{5}{36} = \frac{5}{36}, & \mbox{for \(x \in \{-1, 1\}\)} \\
        \frac{4}{36} = \frac{1}{9}, & \mbox{for \(x \in \{-2, 2\}\)} \\
        \frac{3}{36} = \frac{1}{12}, & \mbox{for \(x \in \{-3, 3\}\)} \\
        \frac{2}{36} = \frac{1}{18}, & \mbox{for \(x \in \{-4, 4\}\)} \\
        \frac{1}{36}, & \mbox{for \(x \in \{-5, 5\}\)} \\
        0, & \mbox{otherwise}
    \end{cases}
\end{equation*}


\section*{Problem 4}

\begin{equation*}
    P_{X}(k) = 
    \begin{cases}
        \frac{1}{4}, & \mbox{for \(k = 1\)} \\
        \frac{1}{8}, & \mbox{for \(k = 2\)} \\
        \frac{1}{8}, & \mbox{for \(k = 3\)} \\
        \frac{1}{2}, & \mbox{for \(k = 4\)} \\
        0, & \mbox{otherwise}
    \end{cases}
\end{equation*}

\begin{equation*}
    P_{Y}(k) = 
    \begin{cases}
        \frac{1}{6}, & \mbox{for \(k = 1\)} \\
        \frac{1}{6}, & \mbox{for \(k = 2\)} \\
        \frac{1}{3}, & \mbox{for \(k = 3\)} \\
        \frac{1}{3}, & \mbox{for \(k = 4\)} \\
        0, & \mbox{otherwise}
    \end{cases}
\end{equation*}

We have:

\subsection*{a.}
\begin{equation*}
    P(X \leq 2 \ \mbox{and} \ Y \leq 2) = P(X \leq 2)P(Y \leq 2) = (\frac{1}{4} + \frac{1}{8}) \cdot (\frac{1}{6} + \frac{1}{6}) =
\end{equation*}
\begin{equation*}
    = \frac{3}{8} \cdot \frac{1}{3} = \frac{1}{8}
\end{equation*}

\subsection*{b.}
\begin{equation*}
    P(X > 2 \ \mbox{or} \ Y > 2) = 1 - P(X \leq 2 \ \mbox{and} \ Y \leq 2) = 1 - \frac{1}{8} = \frac{7}{8}
\end{equation*}

\subsection*{c.}
From the fact that \(X\) and \(Y\) are independent:
\begin{equation*}
    P(X>2|Y>2) = P(X>2) = P(X=3) + P(X=4) = \frac{1}{8} + \frac{1}{4} = \frac{3}{8}
\end{equation*}

\subsection*{d.}
% \begin{equation*}
%     P(X < Y) = P(\{\omega \in \Omega: X(\omega) = 1 \} \cup \{\omega \in \Omega: Y(\omega) \in \{2, 3, 4\} \}) 
% \end{equation*}
\newpage
Using the fact that the sets are disjoint (WHY?):
\begin{equation*}
    \begin{split}
        P(X<Y) & = P\Bigg( \\
        & \Big(\{\omega \in \Omega: X(\omega) = 1 \} \cap \{\omega \in \Omega: Y(\omega) \in \{2, 3, 4\} \} \Big) \\
        & \cup \Big(\{\omega \in \Omega: X(\omega) = 2 \} \cap \{\omega \in \Omega: Y(\omega) \in \{3, 4\} \}\Big) \\
        & \cup \Big(\{\omega \in \Omega: X(\omega) = 3 \} \cap \{\omega \in \Omega: Y(\omega) \in \{4\} \}\Big) \\
        & \Bigg)
    \end{split}
\end{equation*}
\begin{equation*}
    = \frac{1}{4} \cdot (\frac{1}{6} + \frac{1}{3} + \frac{1}{3}) + \frac{1}{8} \cdot (\frac{1}{3} + \frac{1}{3}) + \frac{1}{8} \cdot \frac{1}{3}
\end{equation*}
\begin{equation*}
    = \frac{5}{24} + \frac{1}{12} + \frac{1}{24} = \frac{1}{3}
\end{equation*}


% \begin{multline*}
%     P(X < Y) = P\Bigg( \Big(\{\omega \in \Omega: X(\omega) = 1 \} \cap \{\omega \in \Omega: Y(\omega) \in \{2, 3, 4\} \} \Big) \ \cup \\
%     \Big(\{\omega \in \Omega: X(\omega) = 2 \} \cap \{\omega \in \Omega: Y(\omega) \in \{3, 4\} \}\Big) \ \cup \\
%     \Big(\{\omega \in \Omega: X(\omega) = 3 \} \cap \{\omega \in \Omega: Y(\omega) \in \{4\} \}\Big) \Bigg)
% \end{multline*}

\subsection*{Appendix}

TODO: there are three different ways to attack problem (d.):

\begin{itemize}
    \item via a table X-Y column-row
    \item via the law of total probability
    \item via disjoint events (plus finite additivity of the probability)
\end{itemize}

\section*{Problem 5}
Let \(X\) be a random variable indicating how many students have a car. We want to know what is the probability that more than 30 students
have a car, which will mean, that there were not enough parking spaces for the students. We need to keep choosing
the students from 31 to 50 and in each case multiply that by the chance of having a car vs not having a car. We can use
binomial distribution for that where \(p = 0.5\) is the success (have cars) and \(q = 1 - p\) is the failure (do not have cars).
We are summing up the choices of the students \(\binom{50}{31}, \binom{50}{32}, \binom{50}{33}, etc.\). We have:
\begin{equation*}
    P(\mbox{not enough parking spaces for all the cars}) =  
\end{equation*}
\begin{equation*}
    \sum_{i=31}^{50}\binom{50}{i} \Big(\frac{1}{2}\Big)^{i}(1 - \frac{1}{2})^{50 - i} = 
\end{equation*}
\begin{equation*}
    = \sum_{i=31}^{50}\binom{50}{i} \Big(\frac{1}{2}\Big)^{50} = 0.05946022627971814 \approx 0.05946
\end{equation*}

\section*{Problem 7}

\subsection*{a.}

% \frac{1}{5}\sum_{k=}^{}(\frac{4}{5})^{k-1}

Let \(X \sim Geometric(\frac{1}{5})\):

\begin{equation*}
    P(X > 5) = 1 - P(X \leq 5) = 
\end{equation*}
\begin{equation*}
    = 1 - \sum_{k=1}^{5} \frac{1}{5}(1 - \frac{1}{5})^{k-1} = 1 - \frac{1}{5}\sum_{k=1}^{5}(\frac{4}{5})^{k-1}
\end{equation*}
\begin{equation*}
    \approx 0.32767999999999986 \approx 0.32768
\end{equation*}

\begin{equation*}
    P(2 < X \leq 6) = F_{X}(6) -  F_{X}(2) = \frac{1}{5}\sum_{k=1}^{6}(\frac{4}{5})^{k-1} - \frac{1}{5}\sum_{k=1}^{2}(\frac{4}{5})^{k-1}
\end{equation*}
\begin{equation*}
    = \frac{1}{5}\sum_{k=3}^{6}(\frac{4}{5})^{k-1} \approx 0.3778560000000001 \approx 0.377856
\end{equation*}

\begin{equation*}
    P(X>5|X<8) = \frac{P(5 < X < 8)}{P(X<8)} = \frac{F_{X}(8^{-}) - F_{X}(5)}{F_{X}(8^{-})}
\end{equation*}
\begin{equation*}
    = \frac{\Big(F_{X}(8) - P(X=8)\Big) - F_{X}(5)}{F_{X}(8) - P(X=8)} = 
\end{equation*}
\begin{equation*}
    = \frac{F_{X}(7) - F_{X}(5)}{F_{X}(7)} = 
\end{equation*}
\begin{equation*}
   = \frac{\frac{1}{5}\sum_{k=1}^{7}(\frac{4}{5})^{k-1} - \frac{1}{5}\sum_{k=1}^{5}(\frac{4}{5})^{k-1}}{\frac{1}{5}\sum_{k=1}^{7}(\frac{4}{5})^{k-1}}
\end{equation*}
\begin{equation*}
    \approx \frac{0.7902848 - 0.6723200}{0.7902848} \approx 0.14927
\end{equation*}

\subsection*{b.}

Let \(X \sim Binomial(10, \frac{1}{3})\):


% \sum_{k=0}^{5}\binom{10}{k}\Big(\frac{1}{3}\Big)^{k}\Big(\frac{2}{3}\Big)^{10-k}
\begin{equation*}
    P(X > 5) = 1 - P(X \leq 5) = 1-  \sum_{k=0}^{5}\binom{10}{k}\Big(\frac{1}{3}\Big)^{k}\Big(\frac{2}{3}\Big)^{10-k}
\end{equation*}
\begin{equation*}
    \approx 0.07656353198191379 \approx 0.07656
\end{equation*}
\begin{equation*}
    P(2 < X \leq 6) = F_{X}(6) -  F_{X}(2) = 
\end{equation*}
\begin{equation*}
    = \sum_{k=0}^{6}\binom{10}{k}\Big(\frac{1}{3}\Big)^{k}\Big(\frac{2}{3}\Big)^{10-k} - \sum_{k=0}^{2}\binom{10}{k}\Big(\frac{1}{3}\Big)^{k}\Big(\frac{2}{3}\Big)^{10-k} =
\end{equation*}
\begin{equation*}
    = \sum_{k=3}^{6}\binom{10}{k}\Big(\frac{1}{3}\Big)^{k}\Big(\frac{2}{3}\Big)^{10-k} \approx 0.6811969720062995 \approx 0.68111
\end{equation*}
\begin{equation*}
    P(X>5|X<8) = \frac{\sum_{k=6}^{7}\binom{10}{k}\Big(\frac{1}{3}\Big)^{k}\Big(\frac{2}{3}\Big)^{10-k}}{\sum_{k=0}^{7}\binom{10}{k}\Big(\frac{1}{3}\Big)^{k}\Big(\frac{2}{3}\Big)^{10-k}} =
\end{equation*}
\begin{equation*}
    \approx \frac{0.0731595793324188}{0.996596047350505} \approx 0.0734094616639478 \approx 0.07341
\end{equation*}

\subsection*{c.}

Let \(X \sim Pascal(3, \frac{1}{2})\):

\begin{equation*}
    P(X > 5) = 1 - P(X \leq 5) = 1 - \sum_{k=3}^{5} \binom{k-1}{2}(\frac{1}{2})^{3}(\frac{1}{2})^{k-3}
\end{equation*}
\begin{equation*}
    P(2 < X \leq 6) = \sum_{k=3}^{6} \binom{k-1}{2}(\frac{1}{2})^{3}(\frac{1}{2})^{k-3}
\end{equation*}
\begin{equation*}
    P(X>5|X<8) = \frac{\sum_{k=3}^{7} \binom{k-1}{2}(\frac{1}{2})^{3}(\frac{1}{2})^{k-3} - \sum_{k=3}^{5} \binom{k-1}{2}(\frac{1}{2})^{3}(\frac{1}{2})^{k-3}}{\sum_{k=3}^{7} \binom{k-1}{2}(\frac{1}{2})^{3}(\frac{1}{2})^{k-3}} =
\end{equation*}
\begin{equation*}
    = \frac{\sum_{k=6}^{7} \binom{k-1}{2}(\frac{1}{2})^{3}(\frac{1}{2})^{k-3}}{\sum_{k=3}^{7} \binom{k-1}{2}(\frac{1}{2})^{3}(\frac{1}{2})^{k-3}}
\end{equation*}

\subsection*{d.}

Let \(X \sim Hypergeometric(10, 10, 12)\). Notice that \(R_{X} = \{2, 3, \dots, 10\}\):


\begin{equation*}
    P(X > 5) = 1 - \sum_{x=2}^{5}\frac{\binom{10}{x}\binom{10}{12-x}}{\binom{20}{12}}
\end{equation*}
\begin{equation*}
    P(2 < X \leq 6) = \sum_{x=3}^{6}\frac{\binom{10}{x}\binom{10}{12-x}}{\binom{20}{12}}
\end{equation*}
\begin{equation*}
    P(X>5|X<8) = \frac{\sum_{x=6}^{7}\frac{\binom{10}{x}\binom{10}{12-x}}{\binom{20}{12}}}{\sum_{x=2}^{7}\frac{\binom{10}{x}\binom{10}{12-x}}{\binom{20}{12}}}
\end{equation*}

\subsection*{e.}

Let \(X \sim Poisson(5)\):

\begin{equation*}
    P(X > 5) = 1 - \sum_{k=0}^{5}\frac{e^{- \lambda}\lambda^{k}}{k!}
\end{equation*}
\begin{equation*}
    P(2 < X \leq 6) = \sum_{k=3}^{6}\frac{e^{- \lambda}\lambda^{k}}{k!}
\end{equation*}
\begin{equation*}
    P(X>5|X<8) = \frac{\sum_{k=6}^{7}\frac{e^{- \lambda}\lambda^{k}}{k!}}{\sum_{k=0}^{7}\frac{e^{- \lambda}\lambda^{k}}{k!}}
\end{equation*}


\section*{Problem 8}

\subsection*{a.}

\begin{equation*}
    P(X = 1) = P(S_{1}) = \frac{1}{2}
\end{equation*}

\begin{equation*}
    P(X = 2) = P(\mbox{Failure 1st AND Passed 2nd }) =
\end{equation*}
\begin{equation*}
    = P(F_{1})P(S_{2}) = P(F_{1})(1 - P(F_{2})) = \frac{1}{2} \cdot (1 - \frac{1}{4}) = \frac{3}{8}
\end{equation*}

\begin{equation*}
    P(X = 3) = P(\mbox{Failure 1st AND Failure 2nd AND Success 3rd}) =
\end{equation*}
\begin{equation*}
    P(X = 3) = P(F_{1})P(F_{2})P(S_{3}) = \frac{1}{2} \cdot (\frac{1}{2} \cdot \frac{1}{2}) \cdot (1 - P(F_{3})) =
\end{equation*}
\begin{equation*}
    = \frac{1}{2} \cdot (\frac{1}{2} \cdot \frac{1}{2}) \cdot (1 - (\frac{1}{2} \cdot \frac{1}{2} \cdot \frac{1}{2})) =
\end{equation*}
\begin{equation*}
    = \Big(\frac{1}{2}\Big)^{3} \cdot \Bigg( 1 - \Big(\frac{1}{2}\Big)^{3}\Bigg) = \frac{7}{64}
\end{equation*}

\subsection*{b.}

\begin{equation*}
    P(X=k) = P(F_{1}) \cdot P(F_{2}) \cdot P(F_{3}) \cdot (\dots) \cdot P(S_{k}) = 
\end{equation*}
\begin{equation*}
    \frac{1}{2} \cdot (\frac{1}{2} \cdot \frac{1}{2}) \cdot (\frac{1}{2} \cdot \frac{1}{2} \cdot \frac{1}{2}) \cdot (\dots) \cdot (1 - F_{k}) =
\end{equation*}
\begin{equation*}
    \Big(\frac{1}{2}\Big)^{k-1} \cdot \Bigg(1 - \Big(\frac{1}{2}\Big)^{k}\Bigg) =
\end{equation*}
\begin{equation*}
    \Bigg[\prod_{i=1}^{k-1} \Big(\frac{1}{2}\Big)^{i}\Bigg] \cdot \Bigg(1 - \Big(\frac{1}{2}\Big)^{k}\Bigg) =
\end{equation*}
\begin{equation*}
    \Big(\frac{1}{2}\Big)^{1 + 2 + 3 + \dots + k - 1} \cdot \Bigg(1 - \Big(\frac{1}{2}\Big)^{k}\Bigg) =
\end{equation*}
\begin{equation*}
    \Big(\frac{1}{2}\Big)^{\frac{(k-1)k}{2}} \cdot \Bigg(1 - \Big(\frac{1}{2}\Big)^{k}\Bigg)
\end{equation*}

\subsection*{c.}

\begin{equation*}
    P(X > 2) = 1 - P(X \leq 2) = 1 - (P(X=1) + P(X=2)) = 
\end{equation*}
\begin{equation*}
    1 - \Bigg[\Big(\frac{1}{2}\Big)^{\frac{0}{2}} \cdot \Bigg(1 - \Big(\frac{1}{2}\Big)^{1}\Bigg) + \Big(\frac{1}{2}\Big)^{\frac{2}{2}} \cdot \Bigg(1 - \Big(\frac{1}{2}\Big)^{2}\Bigg)\Bigg] = 0.125
\end{equation*}

\subsection*{d.}

\begin{equation*}
    P(X = 2|X > 1) = \frac{P(X = 2 \ \mbox{and} \ X > 1)}{P(X > 1)} =
\end{equation*}
\begin{equation*}
    \frac{P(X=2)}{P(X > 1)} = \frac{P(X=2)}{1 - P(X \leq 1)} =
\end{equation*}
\begin{equation*}
    \frac{\Big(\frac{1}{2}\Big)^{\frac{2}{2}} \cdot \Bigg(1 - \Big(\frac{1}{2}\Big)^{2}\Bigg)}{1 - \Bigg[\Big(\frac{1}{2}\Big)^{\frac{0}{2}} \cdot \Bigg(1 - \Big(\frac{1}{2}\Big)^{1}\Bigg)\Bigg]} =
\end{equation*}
\begin{equation*}
    \frac{\frac{1}{2} - \Big(\frac{1}{2}\Big)^{3}}{1 - 1 + \frac{1}{2}} = \frac{\frac{3}{8}}{\frac{1}{2}} = \frac{3}{4} = 0.75
\end{equation*}


\section*{Problem 9}
We want to show that:
\begin{equation*}
    P(X > m + l | X > m) = P(X > l), \ \mbox{for} \ m, l \in \{1,2,3, \dots\}
\end{equation*}
We have:

\begin{equation*}
    P(X > m + l | X > m) = \frac{P(X > m + l \ \mbox{and} \  X > m)}{P(X > m)} =
\end{equation*}
\begin{equation*}
    \frac{P(X > m + l > m)}{P(X > m)} = \frac{P(X > m + l)}{P(X > m)} = \frac{1 - P(X \leq m + l) }{1 - P(X \leq m)} = \frac{1 - F_{X}(m + l)}{1 - F_{X}(m)}
\end{equation*}
Recall that the CDF of the Geometric distribution is:

\begin{equation*}
    F_{X}(x) = 1 - (1 - p)^{x} \ \mbox{for \(x \in \{1, 2, 3, \dots\}\)  }
\end{equation*}
\begin{equation*}
    F_{X}(x) = 0 \ \mbox{otherwise}
\end{equation*}

Observe also that:

\begin{equation*}
    1 - (1 - p)^{l} = F(l) = P(X \leq l) = 1 - P(X > l)
\end{equation*}
\begin{equation*}
    1 - (1 - p)^{l} = 1 - P(X > l)
\end{equation*}
\begin{equation*}
    (1 - p)^{l} = P(X > l)
\end{equation*}

Continuing then above, we get:

\begin{equation*}
    P(X > m + l | X > m) = \frac{1 - F_{X}(m + l)}{1 - F_{X}(m)} = \frac{1 - \Big(1 - (1 - p)^{m+l}\Big)}{1 - \Big(1 - (1 - p)^{m}\Big)} =
\end{equation*}
\begin{equation*}
    \frac{(1 - p)^{m+l}}{(1 - p)^{m}} = \frac{(1 - p)^{m}(1 - p)^{l}}{(1 - p)^{m}} = (1 - p)^{l} = P(X > l)
\end{equation*}
\section*{Problem 12}

Let \(X\) be a discrete random variable with the following PMF:

\begin{equation*}
    P_{X}(x) = 
    \begin{cases}
        0.2  & \mbox{for \(x = -2\)} \\
        0.3  & \mbox{for \(x = -1\)} \\
        0.2  & \mbox{for \(x = 0\)} \\
        0.2 & \mbox{for \(x = 1\)} \\
        0.1 & \mbox{for \(x = 2\)} \\
        0 & \mbox{otherwise}
    \end{cases}
\end{equation*}

The CDF is then:


\begin{equation*}
    F_{X}(x) = P(X \leq x) = 
    \begin{cases}
        0  & \mbox{for \(x < -2\)} \\
        \sum_{x \in \{-2\}} P(X = x) & \mbox{for \( -2 \leq x < -1 \)} \\
        \sum_{x \in \{-2, -1\}} P(X = x)  & \mbox{for \( -1 \leq x < 0 \)} \\
        \sum_{x \in \{-2, -1, 0\}} P(X = x) & \mbox{for \( 0 \leq x < 1 \)} \\
        \sum_{x \in \{-2, -1, 0, 1\}} P(X = x)  & \mbox{for \( 1 \leq x < 2 \)} \\
        \sum_{x \in \{-2, -1, 0, 1, 2\}} P(X = x)  & \mbox{for \(2 \leq x\)}
    \end{cases}
\end{equation*}
which is:
\begin{equation*}
    F_{X}(x) = P(X \leq x) = 
    \begin{cases}
        0  & \mbox{for \(x < -2\)} \\
        0.2  & \mbox{for \( -2 \leq x < -1 \)} \\
        0.5  & \mbox{for \( -1 \leq x < 0 \)} \\
        0.7 & \mbox{for \( 0 \leq x < 1 \)} \\
        0.9 & \mbox{for \( 1 \leq x < 2 \)} \\
        1 & \mbox{for \(2 \leq x\)}
    \end{cases}
\end{equation*}

\section*{Problem 13}

The range of the random variable can be found from the jump points of the cdf.
In this case the range is \(R_{X} = \{0, 1, 2, 3\}\). We have then again from the cdf:

\begin{equation*}
    P(X=0) = \frac{1}{6}
\end{equation*}
\begin{equation*}
    P(X=1) = \frac{1}{3}
\end{equation*}
\begin{equation*}
    P(X=2) = \frac{1}{4}
\end{equation*}
\begin{equation*}
    P(X=3) = \frac{1}{4}
\end{equation*}

To summarize:
\begin{equation*}
    P_{X}(X = x) =
    \begin{cases}
        \frac{1}{6}  & \mbox{for \(k = 0\)} \\
        \frac{1}{3}  & \mbox{for \(k = 1\)} \\
        \frac{1}{4}  & \mbox{for \(k \in \{2, 3\}\)} \\
        0 & \mbox{otherwise} 
    \end{cases}
\end{equation*}

\section*{Problem 14}

\subsection*{a.}
\begin{equation*}
    EX = (1 \cdot 0.5) + (2 \cdot 0.3) + (3 \cdot 0.2) = 1.7
\end{equation*}

\subsection*{b.}

\begin{equation*}
    \sigma^{2} = \mathcal{D}^{2}X = \sum_{x_{k} \in R_{X}}(x_{k} - EX)^{2}P(X=x_{k}) = 
\end{equation*}
\begin{equation*}
    \Big((1 - 1.7)^{2} \cdot 0.5 \Big) + \Big((2 - 1.7)^{2} \cdot 0.3 \Big) + \Big((3 - 1.7)^{2} \cdot 0.2 \Big) =
\end{equation*}
\begin{equation*}
    \approx 0.61
\end{equation*}
We could also use (recommended):
\begin{equation*}
    \sigma^{2} = \mathcal{D}^{2}X = E[X^{2}]- [EX]^{2} = \Big((1^{2} \cdot 0.5) + (2^{2} \cdot 0.3) + (3^{2} \cdot 0.2)\Big) - 2.89 = 3.5 - 2.89
\end{equation*}
\begin{equation*}
    \approx 0.6099999999999999 \approx 0.61
\end{equation*}

\begin{equation*}
    \sigma \approx \sqrt{0.61} \approx 0.781025
\end{equation*}

\subsection*{c.}

We use LOTUS with \(g(x) = \frac{2}{x}\):

\begin{equation*}
    EY = E[g(X)] = \sum_{k=1}^{3}\frac{2}{k}P(X=k) = 
\end{equation*}
\begin{equation*}
    \frac{2}{1}\frac{1}{2} + \frac{2}{2}\frac{3}{10} + \frac{2}{3}\frac{2}{10} =
\end{equation*}
\begin{equation*}
    \frac{43}{30} = 1.4(3)
\end{equation*}

\section*{Problem 15}

% Geometric PMF (1 - p)^(k-1)p
Let \(X \sim Geometric(\frac{1}{3})\), and let \(Y = |X - 5|\). Since \(X\) has the Geometric distribution
its range is equal to:

\begin{equation*}
    R_{X} = \{1, 2, 3, 4, 5, \dots\}
\end{equation*}
We shall note that the range of \(Y\) can be written as:

\begin{equation*}
    R_{Y} = \{g(x): \ x \in R_{X}\} = \{|x - 5|: x \in R_{X} \} = 
\end{equation*}
\begin{equation*}
    \{|1-5|, |2-5|, |3-5|, |4-5|, |5-5|, |6-5|, |7-5|, \dots\} =
\end{equation*}
\begin{equation*}
    \{|4|, |3|, |2|, |1|, |0|, |1|, |2|, |3|, |4|, \dots\} = \{0, 1, 2, 3, 4, 5, \dots\}
\end{equation*}

We know the PMF of \(X\) is: 

\begin{equation*}
    P_{X}(X = x) =
    \begin{cases}
        \frac{1}{3} \cdot \Big(\frac{2}{3}\Big)^{x-1}  & \mbox{for} \ x \in \{1, 2, 3, \dots\} \\
        0 & \mbox{otherwise} 
    \end{cases}
\end{equation*}


The PMF of \(Y\) then is:

\begin{equation*}
    P(Y=y) = P(|X - 5| = y) = \sum_{x:g(x) = y}P(X=x)
\end{equation*}


It follows that The PMF of \(Y\) is:

\begin{equation*}
    P_{Y}(Y = y) =
    \begin{cases}
        P(X = 5)  & \mbox{for \(y = 0\)} \\
        P(X = 5 - y) + P(X = 5 + y)  & \mbox{for} \ y \in \{1, 2, 3, 4\} \\
        P(X = 10)  & \mbox{for \(y = 5\)} \\
        P(X = 5 + y)  & \mbox{for} \ y \in \{6, 7, 8, 9, 10, \dots\} \\
        0 & \mbox{otherwise} 
    \end{cases}
\end{equation*}

\begin{equation*}
    P_{Y}(Y = y) =
    \begin{cases}
        \frac{16}{243}  & \mbox{for \(y = 0\)} \\
        \frac{1}{3}\Big(\frac{2}{3}\Big)^{4-y} + \frac{1}{3}\Big(\frac{2}{3}\Big)^{4+y}  & \mbox{for} \ y \in \{1, 2, 3, 4\} \\
        \frac{1}{3}\Big(\frac{2}{3}\Big)^{4+y}  & \mbox{for} \ y \in \{5, 6, 7, 8, 9, 10, \dots\} \\
        0 & \mbox{otherwise} 
    \end{cases}
\end{equation*}

\section*{Problem 16}

The range of \(Y\) is:

\begin{equation*}
    R_{Y} = \{ g(x): x \in R_{X} \} = \{ g(-10), g(-9), g(-8), \cdots, g(0), g(1), \cdots, g(10) \} =
\end{equation*}
\begin{equation*}
    = \{0, 1, 2, 3, 4, 5\}
\end{equation*}

We have for every value in \(R_{Y}\):

\begin{equation*}
    P(Y=0) = \sum_{x:g(x) = y}P(X = x) = \sum_{k \in \{-10, -9, -8, \dots, 0\}}P(X=k) = \frac{11}{21} 
\end{equation*}
\begin{equation*}
    P(Y=y) = \sum_{x:g(x) = y}P(X = x) = P(X=y) = \frac{1}{21} \ \ \text{for} \ \ y \in \{1, 2, 3, 4\}
\end{equation*}
\begin{equation*}
    P(Y=5) = \sum_{x:g(x) = y}P(X = x) = P(X=5) + \sum_{k=6}^{10}P(X=k) = 6 \cdot \frac{1}{21} = \frac{6}{21}
\end{equation*}
\begin{equation*}
    P(Y=y) = 0 \ \ \text{for} \ \ y \notin \{0, 1, 2, 3, 4, 5\}
\end{equation*}

To summarize:

\begin{equation*}
    P(Y = y) =
    \begin{cases}
        \frac{11}{21} & \mbox{for} \ y = 0 \\
        \frac{1}{21}  & \mbox{for} \ y \in \{1, 2, 3, 4\} \\
        \frac{6}{21}  & \mbox{for} \ y = 5 \\
        0 & \mbox{otherwise} 
    \end{cases}
\end{equation*}


\section*{Problem 17}

Let \(X \sim Geometric(p)\). We seek to calculate \(Var(X)\). We can use the following formula:
\begin{equation*}
    Var(X) = E[X^{2}] - [EX]^{2}
\end{equation*}
We will start with \(EX\). Recall the formula for an expected value:

\begin{equation*}
    EX = \sum_{x_{k} \in R_{X}} x_{k}P(X = x_{k}) = \sum_{x_{k} \in R_{X}}x_{k}P_{X}(x_k)
\end{equation*}
Since \(X\) has a geometric distribution, the PMF of \(X\) is:

\begin{equation*}
    P_{X}(k) =
    \begin{cases}
        p(1 - p)^{k-1} & \mbox{for} \ k = 1, 2, 3, \dots \\
        0 & \mbox{otherwise} 
    \end{cases}
\end{equation*}
where \(0 < p < 1\). We have:

\begin{equation*}
    EX = \sum_{k = 1}^{\infty}kP_{X}(x_k) = \sum_{k = 1}^{\infty} kp(1 - p)^{k-1} = p \sum_{k = 1}^{\infty} k(1 - p)^{k-1}
\end{equation*}
\begin{equation*}
    = p \cdot \Big( 1(1 - p)^{0} + 2(1 - p)^{1} + 3(1 - p)^{2} + \dots \Big)
\end{equation*}

Observe that we can use the geometric series to calculate the sum at hand. It is fairly known that:
\begin{equation*}
    \sum_{k=0}^{\infty}px^{k} = \sum_{k=1}^{\infty}px^{k-1}  = \frac{p}{1 - x} \ \mbox{for} \ |x| < 1, p \neq 0
\end{equation*}


We can calculate the derivative of the sum:

\begin{equation*}
    \frac{d}{dx}\Big(\frac{p}{1 - x}\Big) = \frac{p}{(1 - x)^{2}}
\end{equation*}

Let \(x = 1 - p\). Notice that \(|x| < 1\). Then we see:


\begin{equation*}
    p \sum_{k = 1}^{\infty} k(1 - p)^{k-1} = p \sum_{k = 1}^{\infty} kx^{k-1} = 
\end{equation*}
\begin{equation*}
    = p \cdot \Big( 1 + 2x + 3x^{2} + 4x^{3} + \dots \Big) 
\end{equation*}
\begin{equation*}
    = p \cdot \Big( \frac{dx}{dx} + \frac{dx^{2}}{dx} + \frac{dx^{3}}{dx} + \frac{dx^{4}}{dx} + \dots \Big)
\end{equation*}
\begin{equation*}
    = p \cdot \frac{d}{dx}\Big( x + x^{2} + x^{3} + x^{4} + \dots  \Big)
\end{equation*}
\begin{equation*}
    = p \cdot \frac{d}{dx}\Big( \sum_{k=1}^{\infty} x^{k}  \Big) = \frac{d}{dx} \frac{p}{1-x} = \frac{p}{(1 - x)^{2}} = \frac{p}{p^{2}} = \frac{1}{p}
\end{equation*}

Then:

\begin{equation*}
    [EX]^{2} = \frac{1}{p^{2}}
\end{equation*}

Next, we calculate \(EX^{2}\). From the definition of the expected value:

\begin{equation*}
    EX^{2} =  \sum_{x_{k} \in R_{X}} x_{k}^{2} P(X = x_{k}) = \sum_{x_{k} \in R_{X}}x_{k}^{2}P_{X}(x_k) =
\end{equation*}
\begin{equation*}
    = p \sum_{k=1}^{\infty} k^{2}(1 - p)^{k-1} = p \sum_{k=1}^{\infty} k^{2}x^{k-1}
\end{equation*}
with \(x = 1 - p\). We can differentiate the geometric series twice to discover a helpful fact for our case:

\begin{equation*}
    \sum_{k=0}^{\infty} x^{k} = \frac{1}{1-x}
\end{equation*}

\begin{equation*}
    \Big(\sum_{k=0}^{\infty} x^{k} \Big)^{'} = \sum_{k=0}^{\infty} (x^{k})^{'} = (x^{0})^{'} + (x^{1})^{'} + (x^{2})^{'} + (x^{3})^{'} + \dots
\end{equation*}
\begin{equation*}
    = 0 + 1 + 2x + 3x^{2} + \dots = 1 + 2x + 3x^{2} + \dots
\end{equation*}
\begin{equation*}
    \sum_{k=1}^{\infty} kx^{k-1} = \frac{1}{(x-1)^{2}}
\end{equation*}

\begin{equation*}
    \Big(\sum_{k=1}^{\infty} kx^{k-1}\Big)^{'} = \sum_{k=1}^{\infty} (kx^{k-1})^{'} = (1x^{0})^{'} + (2x^{1})^{'} + (3x^{2})^{'} + (4x^{3})^{'} + 
\end{equation*}
\begin{equation*}
    + \  (5x^{4})^{'} + (6x^{5})^{'} + (7x^{6})^{'} + \dots
\end{equation*}
\begin{equation*}
    = 0 + 2 + 6x + 12x^{2} + 20x^{3} + 30x^{4} + 42x^{5} + \dots
\end{equation*}
\begin{equation*} 
    = (2 \cdot 1) + (3 \cdot 2)x + (4 \cdot 3)x^{2} + (5 \cdot 4)x^{3} + (6 \cdot 5)x^{4} + (7 \cdot 6)x^{5} + \dots
\end{equation*}
\begin{equation*}
    = \sum_{k=2}^{\infty} k(k-1)x^{k-2} = - \frac{2}{(x-1)^{3}} =  \frac{2}{(1-x)^{3}}
\end{equation*}

We use the last line to do the following:
\begin{equation*}
    E[X(X - 1)]
\end{equation*}



\section*{Problem 19}


Suppose that \(Y = -2X + 3\). Notice that \(X = \frac{3 - Y}{2}\). We know that \(EY = 1\) and \(EY^{2} = 9\). We have:

\begin{equation*}
    EX = E[\frac{3 - Y}{2}] = \frac{3}{2} - \frac{EY}{2} = \frac{3}{2} - \frac{1}{2} = 1
\end{equation*}
\begin{equation*}
    Var(X) = E[X^{2}] - [EX]^{2} = E[X^{2}] - 1 = E[\frac{9 - 6Y + Y^{2}}{4}] - 1 = 
\end{equation*}
\begin{equation*}
    = \frac{9}{4} - \frac{6}{4}EY + \frac{1}{4}E[Y^{2}] - 1 = \frac{9}{4} - \frac{6}{4} + \frac{9}{4} - 1 = 2
\end{equation*}

\section*{Problem 20}

\subsection*{(a)}

The range of the random variable is \(R_{X} = \{1, 2, 3, 4, 5, 6\}\). The PMF of \(X\) is:

\begin{equation*}
    P(X = k) =
    \begin{cases}
        \frac{100}{1000} & \mbox{for} \ k \in \{1, 5, 6\} \\
        \frac{200}{1000} & \mbox{for} \ k \in \{2, 4\} \\
        \frac{300}{1000} & \mbox{for} \ k = 3 \\
        0 & \mbox{otherwise} 
    \end{cases}
\end{equation*}

And the expected value is:

\begin{multline*}
    EX = \sum_{k=1}^{6} kP(X=k) = \frac{1 \cdot 100}{1000} + \frac{2 \cdot 200}{1000}  \\ 
    + \frac{3 \cdot 300}{1000} + \frac{4 \cdot 200}{1000} + \frac{5 \cdot 100}{1000} + \frac{6 \cdot 100}{1000} =
\end{multline*}
\begin{equation*}
    = \frac{100 + 400 + 900 + 800 + 500 + 600}{1000} = \frac{3300}{1000} = \frac{33}{10} = 3.3
\end{equation*}

\subsection*{(b)}

The range of the random variable is \(R_{Y} = \{1, 2, 3, 4, 5, 6\}\). The PMF of \(Y\) is:

\begin{equation*}
    P(Y = k) =
    \begin{cases}
        \frac{100}{3300} & \mbox{for} \ k = 1 \\
        \frac{400}{3300} & \mbox{for} \ k = 2 \\
        \frac{900}{3300} & \mbox{for} \ k = 3 \\
        \frac{800}{3300} & \mbox{for} \ k = 4 \\
        \frac{500}{3300} & \mbox{for} \ k = 5 \\
        \frac{600}{3300} & \mbox{for} \ k = 6 \\
        0 & \mbox{otherwise} 
    \end{cases}
\end{equation*}

And the expected value is:

\begin{multline*}
    EY = \sum_{k=1}^{6} kP(Y=k) = \frac{1 \cdot 100}{3300} + \frac{2 \cdot 400}{3300}  \\ 
    + \frac{3 \cdot 900}{3300} + \frac{4 \cdot 800}{3300} + \frac{5 \cdot 500}{3300} + \frac{6 \cdot 600}{3300} =
\end{multline*}
\begin{equation*}
    = \frac{100 + 800 + 2700 + 3200 + 2500 + 3600}{3300} = \frac{12900}{3300} = \frac{129}{33} = \frac{43}{11} \approx 3.91 
\end{equation*}



\section*{Problem 21}

\subsection*{a.}


Because there are \(N\) different types of coupons, we can consider \(N\) trials for every
coupon type using the geometric distribution. So for the first coupon, we create a geometric
distribution, then for the second type, then for the third type and so on and so on.

For every trial (sequences of coupons) of a unique coupon we can show its success and failures
(\(p\) - success because we took our coupon of interest, \(q\) - failure, because it was a different one).
For example:

\begin{multline*}
     \\
    p \\ 
    qp \\ 
    qqp \\ 
    qqqp \\ 
    qqqqp \\ 
    qqqqp \\ 
\end{multline*}
and so on and so on. Using now the language of random variables, we can consider adding the results
of \(N-i\) components (\(i = 0, 1, 2, 3, \dots, N-1\)) like this:

\begin{equation*}
    X = X_{0} + X_{1} + X_{2} + \dots + X_{N-1}
\end{equation*}

with \(X\) being the total number of throws needed to see every coupon at least once
. The first throw will always yield a unique coupon (first occurence) with probability \(\frac{N}{N} = 1\).
Next, we can remove this coupon and focus on the \(N-1\) remaining ones so we wait for the success with the probability
of \(\frac{N-1}{N}\). Next, we can skip the two of the coupons we have already seen and wait for the next one
with the probability of success equal to \(\frac{N-2}{N}\) etc. We conclude that every random variable has
distribution:

\begin{equation*}
    X_{i} \sim Geo(\frac{N-i}{N}) \ \ \mbox{for i = 0, 1, 2, 3, \(\dots\) , \(N-1\)}
\end{equation*}

\subsection*{b.}

We use the linearity of expected value. We have:

\begin{equation*}
    EX = E[X_{0} + X_{1} + \dots + X_{N-1}] = EX_{0} + E_{1} + \dots + EX_{N-1}
\end{equation*}

Because every component (a random variable) has geometric distribution, we sum up the components together
remembering that \(EX_{i} = \frac{N}{N - i} \):

\begin{equation*}
    EX = \frac{N}{N} + \frac{N}{N-1} + \frac{N}{N-2} + \frac{N}{N-3} + \frac{N}{N-4} + \dots + \frac{N}{1}
\end{equation*}
\begin{equation*}
    =  \frac{N}{1} + \frac{N}{2} + \frac{N}{3} + \frac{N}{4} + \frac{N}{5} + \dots + \frac{N}{N}
\end{equation*}
\begin{equation*}
    = N \cdot \Big( \frac{1}{1} + \frac{1}{2} + \frac{1}{3} + \dots + \frac{1}{N} \Big)
\end{equation*}
\begin{equation*}
    = N \sum_{i=1}^{N} \frac{1}{i}
\end{equation*}

\section*{Problem 22}

\subsection*{a.}

We set \(X\) as a random variable indicating the amount of money that the player wins.
The coin tosses are independent so we can model the game using geometric distribution with
the success \(p = \frac{1}{2}\). The range is equal to \(R_{X} = \{2^{0}, 2^{1}, 2^{2}, 2^{3}, \dots \}\).
We proceed to calculate \(EX\):

\begin{equation*}
    EX = \sum_{k=1}^{\infty} \Big(2^{k-1}\Big) \cdot \frac{1}{2} \cdot \Big(\frac{1}{2}\Big)^{k-1} =
\end{equation*}
\begin{equation*}
    = \frac{1}{2} \sum_{k=1}^{\infty} \Big(2^{k-1}\Big) \Big(\frac{1}{2}\Big)^{k-1} =
\end{equation*}
\begin{equation*}
    = \frac{1}{2} \sum_{k=1}^{\infty} 1^{k-1}
\end{equation*}
The series diverges so \(EX\) \(\rightarrow\) \(\infty\).


\subsection*{b.}


\begin{equation*}
    P(X > 65) = 1 - P(X \leq 65) = 1 - \sum_{k=1}^{7} P_{X}(k) = 1 - \sum_{k=1}^{7} \Big(\frac{1}{2}\Big)^{k} =
\end{equation*}
\begin{equation*}
    = \frac{1}{128} = 0.0078125
\end{equation*}

\subsection*{c.}

Let \(Y\) be a random variable representing the money earned with the constraint that the casino will pay \(2^{30}\) max.
Then:

\begin{equation*}
    R_{Y} =
    \begin{cases}
        2^{k-1} & \mbox{for} \ k = 1, \dots, 30, k \in \mathbb{N} \\
        2^{30} & \mbox{for} \ k > 30, k \in \mathbb{N}\\
    \end{cases}
\end{equation*}

\begin{equation*}
    EY = \Big( \frac{1}{2} \sum_{k=1}^{30} 1^{k-1} \Big) + \sum_{k=31}^{\infty} 2^{30} \cdot \big(\frac{1}{2}\big) \big( \frac{1}{2} \big)^{k-1} = 
\end{equation*}
\begin{equation*}
    = 15 + \sum_{k=31}^{\infty} 2^{30} \cdot \big( \frac{1}{2} \big)^{k}
\end{equation*}
\begin{equation*}
    = 15 + \sum_{k=31}^{\infty} 2^{30} \cdot \big( \frac{1}{2} \big)^{k}
\end{equation*}
\begin{equation*}
    = 15 + 1 = 16
\end{equation*}


\section*{Problem 23}

Let \(X\) be a random variable with mean \(EX = \mu \). Define the function \(f(\alpha)\) as:

\begin{equation*}
    f(\alpha) = E[(X - \alpha)^{2}] = E[X^{2} - 2X\alpha + \alpha^{2}]
\end{equation*}

To find the value that minimizes \(f\) we should calculate the derivative of \(f\) and then compare it to \(0\):

\begin{equation*}
    f'(\alpha) = \big( EX^{2} - 2\alpha EX + E[\alpha^{2}]\big)' = -2EX + 2 \alpha = 2(\alpha - EX) = 2(\alpha - \mu)
\end{equation*}

\begin{equation*}
    0 = 2(\alpha - \mu) \implies \mu = EX = \alpha
\end{equation*}

So the value that minimizes \(f\) is \(EX\).

\subsection*{Note to myself}

We get \(Var(X)\) when we put the value that minimizes \(f\) as an argument:

\begin{equation*}
    f(\mu) = EX^{2} - 2 \mu \mu + \mu^{2} = EX^{2} - 2 \mu^{2} + \mu^{2} =  EX^{2} - \mu^{2} 
\end{equation*}


\section*{Problem 24}

\subsection*{Using Law of total expectation}

To maximize \(EW\) one must carefully decide on rolling vs re-rolling. If on average, we know
that the dice roll is \(3.5\), we must re-roll only when we have \(1, 2, 3\). Otherwise we
do not roll again. Let \(X\) indicate not-rolling and \(Y\) rolling again.
Observe that \(P(X) = P(Y) = \frac{1}{2}\).
Then:

\begin{equation*}
    EW = E[W|X]P(X) + E[W|Y]P(Y) = 
\end{equation*}
\begin{equation*}
    = \Big((4 + 5 + 6) \cdot \frac{1}{3}\Big)\frac{1}{2} + \Big((1 + 2 + 3 + 4 + 5 + 6) \cdot \frac{1}{6}\Big)\frac{1}{2} =
\end{equation*}
\begin{equation*}
    = 5 \cdot \frac{1}{2} + 3.5 \cdot \frac{1}{2} = 2.5 + 1.75 = 4.25
\end{equation*}

\subsection*{Reframing the problem using \(\max\) function}


Consider a function representing your maximum prize from entering the game: \(W(k) = \max(k, 3.5)\) for \(k = 1, 2, \dots, 6\):

\begin{equation*}
    W(k) = \max(k, 3.5) =
    \begin{cases}
        3.5 \ \mbox{for} \  k \in \{1, 2, 3\} \\
        k \ \mbox{for} \  k \in \{4, 5, 6\} \\
    \end{cases}
\end{equation*}

If we calculate \(EW\) we get:

\begin{equation*}
    EW = \frac{1}{6} \cdot 3.5 + \frac{1}{6} \cdot 3.5 + \frac{1}{6} \cdot 3.5 + \frac{1}{6} \cdot 4 + \frac{1}{6} \cdot 5 + \frac{1}{6} \cdot 6 =
\end{equation*}
\begin{equation*}
    = \frac{1}{6}(3.5 + 3.5 + 3.5) + \frac{1}{6}(4 + 5 + 6) = \frac{10.5}{6} + \frac{15}{6} = \frac{25.5}{6} = 4.25
\end{equation*}

which is the same score we get in method above.


\section*{Problem 25}

Observe that:

\begin{equation*}
    P(X \geq m) \geq \frac{1}{2} \iff 1 - P(X < m) \geq \frac{1}{2} \iff P(X < m) \leq \frac{1}{2} \iff 
\end{equation*}
\begin{equation*}
    \iff F_{X}(m) - P_{X}(m) \leq \frac{1}{2} \label{median_first_case} \tag{44}
\end{equation*}

As for the second case:
\begin{equation*}
    F_{X}(m) = P(X \leq m) \geq \frac{1}{2} \label{median_second_case} \tag{69}
\end{equation*}

\subsection*{a.}

First, we shall calculate the CDF of \(X\):

\begin{equation*}
    F_{X}(m) = 
    \begin{cases}
        0 \ \mbox{for} \ m \in  (- \infty, 1) \\
        0.4 \ \mbox{for} \ m \in  [1, 2) \\
        0.7 \ \mbox{for} \ m \in  [2, 3) \\
        1 \ \mbox{for} \ m \in [3, + \infty) \\
    \end{cases}
\end{equation*}

In order for \eqref{median_second_case} to happen, \(m \geq 2\). Suppose \(m = 2\). Then:

\begin{equation*}
    F_{X}(2) - P_{X}(2) = 0.7 - 0.3 = 0.4 \leq \frac{1}{2}
\end{equation*}

Suppose now that \(m = 3\). We have:

\begin{equation*}
    F_{X}(3) - P_{X}(3) = 1 - 0.3 = 0.7 > \frac{1}{2}
\end{equation*}
and so we must exclude \(m = 3\) (and simultaneously cases where \(P_{X}\) is zero).
\\
We conclude that the median of \(X\) is \(m = 2\) (unique in this case).


\subsection*{b.}

In order for \eqref{median_second_case} to happen, \(m \geq 3\). We consider now the remaining options for \eqref{median_first_case}:
\begin{equation*}
    F(m) - P(m) = \frac{3}{6} - \frac{1}{6} = \frac{2}{6} \ \ \mbox{for} \ \ m = 3 
\end{equation*}
\begin{equation*}
    F(m) - P(m) = \frac{3}{6} - 0 = \frac{3}{6} \ \ \mbox{for} \ \ m \in (3, 4)
\end{equation*}
\begin{equation*}
    F(m) - P(m) = \frac{4}{6} - \frac{1}{6} = \frac{3}{6} \ \ m = 4
\end{equation*}
\begin{equation*}
    F(m) - P(m) = \frac{4}{6} - 0 = \frac{4}{6} \ \ m \in (4, 5)
\end{equation*}
\begin{equation*}
    F(m) - P(m) = \frac{5}{6} - \frac{1}{6} = \frac{4}{6} \ \ m = 5
\end{equation*}
\begin{equation*}
    F(m) - P(m) = \frac{5}{6} - 0 = \frac{5}{6} \ \ m \in (5, 6)
\end{equation*}
\begin{equation*}
    F(m) - P(m) = \frac{6}{6} - \frac{1}{6} = \frac{5}{6} \ \ m = 6
\end{equation*}
\begin{equation*}
    F(m) - P(m) = \frac{6}{6} - 0 = 1 \ \ \mbox{for} \ \ m \in (6, + \infty)
\end{equation*}

We take probabilities less than or equal to \(\frac{1}{2}\) and so we conclude that the median for rolling a fair die
is \(m \in [3, 4]\) and it is not unique.

\subsection*{c.}

Assume \(X \sim Geometric(p)\) with \(0 < p < 1\). Let \(q = 1 - p\). Next, calculate probabilities:

\begin{equation*}
    P(X \geq m) = \sum_{k=\lceil m \rceil}^{\infty} q^{k-1}p = p \sum_{k=\lceil m \rceil}^{\infty} q^{k-1} =
\end{equation*}
\begin{equation*}
    = pq^{\lceil m \rceil - 1}\Big( q^{0} + q^{1} + q^{2} +  q^{3} \dots \Big) = pq^{\lceil m \rceil - 1} \sum_{k=0}^{\infty} q^{k}
\end{equation*}

Using the sum of an infinite geometric series we get:
\begin{equation*}
    pq^{\lceil m \rceil - 1} \sum_{k=0}^{\infty} q^{k} = pq^{\lceil m \rceil - 1}\frac{1}{1 - q} = 
\end{equation*}
\begin{equation*}
    pq^{\lceil m \rceil - 1}\frac{1}{p} = q^{\lceil m \rceil - 1}
\end{equation*}

As for the second probability term we get:
\begin{equation*}
    P(X \leq x) = \sum_{k=1}^{\lfloor m \rfloor} q^{k-1}p = p \sum_{k=1}^{\lfloor m \rfloor} q^{k-1} =
\end{equation*}
\begin{equation*}
    = p\Big( q^{0} + q^{1} + q^{2} + \dots + q^{\lfloor m \rfloor - 1} \Big) = p \sum_{k=0}^{\lfloor m \rfloor - 1} q^{k}
\end{equation*}

From the sum of a finite geometric series:

\begin{equation*}
    p \sum_{k=0}^{\lfloor m \rfloor - 1} q^{k} = p \frac{1 - q^{\lfloor m \rfloor - 1 + 1}}{1 - q} = 1 - q^{\lfloor m \rfloor}
\end{equation*}


We need to find \(m\) based on two conditions:

\begin{equation*}
    P(X \geq m) = q^{\lceil m \rceil - 1} \geq \frac{1}{2} \label{median_first_geo} \tag{*}
\end{equation*}
\begin{equation*}
    P(X \leq m) = 1 - q^{\lfloor m \rfloor} \geq \frac{1}{2} \implies \frac{1}{2} \geq q^{\lfloor m \rfloor} \label{median_second_geo} \tag{**}
\end{equation*}

First notice that one can make the following transformation using the change of base formula for logarithms:

\begin{equation*}
    \log_{q}\frac{1}{2} = \frac{\log_{2}\frac{1}{2}}{\log_{2}q} = \frac{\log_{2}2^{-1}}{\log_{2}q} = - \frac{\log_{2}2}{\log_{2}q} = \frac{1}{- \log_{2}q} = \frac{1}{\log_{2}q^{-1}} = \frac{1}{\log_{2}\frac{1}{q}}
\end{equation*}
For \eqref{median_first_geo} we make the following transformations:

\begin{equation*}
    q^{\lceil m \rceil - 1} \geq \frac{1}{2}
\end{equation*}
\begin{equation*}
    \log_{q}q^{\lceil m \rceil - 1} \geq \log_{q}\frac{1}{2} 
\end{equation*}
\begin{equation*}
    \lceil m \rceil - 1 \geq \frac{1}{\log_{2}\frac{1}{q}}
\end{equation*}
\begin{equation*}
    \lceil m \rceil \geq \frac{1}{\log_{2}\frac{1}{q}} + 1
\end{equation*}

And now for \eqref{median_second_geo}:

\begin{equation*}
    \frac{1}{2} \geq q^{\lfloor m \rfloor}
\end{equation*}
\begin{equation*}
    \log_{q}\frac{1}{2} \geq \log_{q}q^{\lfloor m \rfloor}
\end{equation*}
\begin{equation*}
    \frac{1}{\log_{2}\frac{1}{q}} \geq q^{\lfloor m \rfloor}
\end{equation*}

\end{document}